% LuaLaTeX Document - Generated by Media Scribe Workflow
\documentclass[a4paper,10pt,twocolumn]{ltjsarticle}

% LuaLaTeX用フォント設定パッケージ
\usepackage{luatexja-fontspec}
\usepackage{amsmath,amssymb}
\usepackage{unicode-math}

% ====================
% 欧文フォント設定 (Libertinus)
% ====================
\setmainfont{Libertinus Serif}[
    BoldFont = {Libertinus Serif Bold},
    ItalicFont = {Libertinus Serif Italic},
    BoldItalicFont = {Libertinus Serif Bold Italic}
]
\setsansfont{Libertinus Sans}[
    BoldFont = {Libertinus Sans Bold},
    ItalicFont = {Libertinus Sans Italic}
]
\setmonofont{DejaVu Sans Mono}[Scale=0.9]

% ====================
% 日本語フォント設定 (原ノ味フォント)
% ====================
\setmainjfont{HaranoAjiMincho-Regular}[
    BoldFont = {HaranoAjiGothic-Medium},
    ItalicFont = {HaranoAjiMincho-Regular},
    BoldItalicFont = {HaranoAjiGothic-Bold}
]
\setsansjfont{HaranoAjiGothic-Regular}[
    BoldFont = {HaranoAjiGothic-Bold}
]
\setmonojfont{HaranoAjiGothic-Regular}

% ====================
% 数式フォント設定 (Libertinus Math)
% ====================
\setmathfont{Libertinus Math}

% ファイル生成日時（JST）
\newcommand{\generatedDate}{2026-01-12}
\newcommand{\generatedTime}{16:58}

% その他のパッケージ
\usepackage[margin=20mm]{geometry}
\usepackage{hyperref}
\usepackage{booktabs}
\usepackage{array}
\usepackage{tabularx}
\usepackage{enumitem}
\usepackage{graphicx}
\usepackage{ascmac}

% ハイパーリンクの色設定
\hypersetup{
    colorlinks=true,
    linkcolor=blue,
    urlcolor=blue,
    citecolor=blue
}

% ヘッダー・フッター設定
\usepackage{fancyhdr}
\usepackage{lastpage}
\pagestyle{fancy}
\fancyhf{}
\fancyhead[R]{\small \generatedDate\ \generatedTime\ JST (\thepage/\pageref{LastPage})}
\renewcommand{\headrulewidth}{0.4pt}

% 1ページ目のスタイル
\fancypagestyle{firstpage}{
    \fancyhf{}
    \fancyhead[R]{\small \generatedDate\ \generatedTime\ JST}
    \renewcommand{\headrulewidth}{0.4pt}
}

% Y列タイプ定義（tabularx用）
\newcolumntype{Y}{>{\raggedright\arraybackslash}X}

\title{\textbf{Media Scribe Workflow\\設定ファイル設計とリファクタリング}}
\author{massy-Claude Dialogue}
\date{}

\begin{document}
\maketitle
\thispagestyle{firstpage}

\tableofcontents

% ============================================================
\section{概要}
% ============================================================

本文書は、Media Scribe Workflow（MSW）の設定ファイル設計とリファクタリング作業の記録である。VCE（Video Chapter Editor）を中心としたメディア処理ワークフローにおいて、テンプレートとプロジェクト設定の分離、編集方針の明確化、LaTeX設定の共通化などを実施した。

\subsection{背景}

MSWは、動画・音声メディアから字幕（SRT）を取得し、LuaTeX形式のレポートを生成するワークフローツール群である。オーケストラ・リハーサル、国際会議、ヨガレッスン、学術講義、楽器レッスンなど、多様な用途に対応する必要があった。

\subsection{課題}

\begin{itemize}[noitemsep]
    \item テンプレート間でのLaTeX設定の重複
    \item \texttt{fidelity}設定の意味が分野で異なる
    \item 静的な用語集の保守限界
    \item 特定講師への依存（horn\_lesson）
\end{itemize}

% ============================================================
\section{設計原則}
% ============================================================

\subsection{三層構造}

設定ファイルを以下の三層に分離した。

\vspace{0.5\baselineskip}
\noindent{\footnotesize
\begin{tabularx}{\linewidth}{@{}lX@{}}
\toprule
層 & 責務 \\
\midrule
defaults.yaml & グローバル設定（著者情報、パス、編集方針デフォルト） \\
templates/*.yaml & 分野固有の構造・セクション定義（差分のみ） \\
project.msw.json & プロジェクト固有データ \\
\bottomrule
\end{tabularx}
}
\vspace{0.5\baselineskip}

\subsection{ファイル形式の選択}

\begin{itemize}[noitemsep]
    \item \textbf{YAML}: 人間が編集するファイル（defaults, templates）
    \item \textbf{JSON}: 機械が管理するファイル（project.msw.json）
\end{itemize}

% ============================================================
\section{editorial\_policyの設計}
% ============================================================

従来の\texttt{fidelity}設定を、より明確な\texttt{editorial\_policy}に再設計した。

\subsection{summarization（要約の許可レベル）}

\vspace{0.5\baselineskip}
\noindent{\footnotesize
\begin{tabularx}{\linewidth}{@{}lX@{}}
\toprule
値 & 説明 \\
\midrule
none & 原文を忠実に保持（会議記録等） \\
minimal & 冗長部分のみ削除可 \\
allowed & 適度な要約を許可 \\
\bottomrule
\end{tabularx}
}
\vspace{0.5\baselineskip}

\subsection{supplementation（補足の積極性）}

\vspace{0.5\baselineskip}
\noindent{\footnotesize
\begin{tabularx}{\linewidth}{@{}lX@{}}
\toprule
値 & 説明 \\
\midrule
minimal & ファイル内の情報のみ使用 \\
contextual & 文脈から推測される補足を許可 \\
aggressive & 一般知識に基づく積極的な補足（初心者向け） \\
\bottomrule
\end{tabularx}
}
\vspace{0.5\baselineskip}

\subsection{target\_audience（対象読者）}

\vspace{0.5\baselineskip}
\noindent{\footnotesize
\begin{tabularx}{\linewidth}{@{}lX@{}}
\toprule
値 & 説明 \\
\midrule
expert & 専門家向け（用語解説最小限） \\
intermediate & 中級者向け \\
beginner & 初心者向け（用語解説充実） \\
\bottomrule
\end{tabularx}
}
\vspace{0.5\baselineskip}

\subsection{テンプレート別のeditorial\_policy}

\vspace{0.5\baselineskip}
\noindent{\footnotesize
\begin{tabularx}{\linewidth}{@{}lYYY@{}}
\toprule
テンプレート & summarization & supplementation & target\_audience \\
\midrule
rehearsal & none & contextual & intermediate \\
meeting & none & minimal & expert \\
yoga & minimal & aggressive & beginner \\
lecture & none & contextual & intermediate \\
instrument\_lesson & minimal & aggressive & beginner \\
\bottomrule
\end{tabularx}
}
\vspace{0.5\baselineskip}

% ============================================================
\section{用語集の動的生成方式}
% ============================================================

\begin{itembox}[l]{問い}
用語などのドメイン知識は、プロジェクトで持つには限界があると思いますので、最低限の固有知識に止め、外部の調査に基づく方が合理的で品質を高く維持できるのではないかと思いますがいかがでしょう。
\end{itembox}

おっしゃる通りです。

\subsection{静的管理の限界}

\begin{itemize}[noitemsep]
    \item 網羅性の限界：どんなに用語を追加しても足りない
    \item 陳腐化：情報が古くなる
    \item 保守コスト：更新作業が負担に
    \item 品質のばらつき：手作業で追加すると精度が不安定
\end{itemize}

\subsection{動的生成方式}

\begin{verbatim}
terminology:
  generation:
    method: "external_research"
    citation_required: true
  static_terms: []
  hint_categories:
    - "速度記号"
    - "強弱記号"
\end{verbatim}

\texttt{hint\_categories}は生成時の参考情報として機能し、実際の用語解説はClaude + 外部調査で動的に生成する。

% ============================================================
\section{LaTeX設定の共通化}
% ============================================================

\begin{itembox}[l]{問い}
luatexの設定も共通にしたいので、切り出してください。最近最も気に入っているのは\texttt{\~{}/.claude/commands/luatex\_dialog.md}です。
\end{itembox}

\texttt{luatex\_dialog.md}を基準とし、\texttt{luatex-settings.yaml}として切り出した。

\subsection{共通設定項目}

\begin{itemize}[noitemsep]
    \item ドキュメントクラス：\texttt{ltjsarticle}、2段組
    \item 欧文フォント：Libertinus Serif/Sans/Math
    \item 日本語フォント：原ノ味明朝/ゴシック
    \item 等幅フォント：DejaVu Sans Mono
    \item ハイパーリンク：青色
    \item ヘッダー：生成日時 JST（ページ/全ページ）
    \item 表：縦線なし、booktabs使用
\end{itemize}

\subsection{ファイル構成}

\begin{verbatim}
~/.config/msw/
├── defaults.yaml
├── luatex-settings.yaml
└── templates/
    ├── rehearsal.yaml
    ├── meeting.yaml
    ├── yoga.yaml
    ├── lecture.yaml
    └── instrument_lesson.yaml
\end{verbatim}

% ============================================================
\section{テンプレートの簡素化}
% ============================================================

各テンプレートを差分のみ定義する形式にリファクタリングした。

\subsection{Before/After比較}

\vspace{0.5\baselineskip}
\noindent{\footnotesize
\begin{tabularx}{\linewidth}{@{}lYY@{}}
\toprule
テンプレート & Before（行数） & After（行数） \\
\midrule
rehearsal.yaml & 154 & 87 \\
meeting.yaml & 181 & 90 \\
yoga.yaml & 171 & 83 \\
lecture.yaml & 174 & 119 \\
horn\_lesson.yaml → instrument\_lesson.yaml & 95 & 82 \\
\bottomrule
\end{tabularx}
}
\vspace{0.5\baselineskip}

\subsection{差分定義の例}

\begin{verbatim}
template_name: "yoga"
extends: "defaults"

editorial_policy:
  summarization: "minimal"
  supplementation: "aggressive"
  target_audience: "beginner"

timestamp:
  format: "[HH:MM:SS.mmm]"

report:
  sections:
    - id: overview
      title: "レッスン概要"
    # ... 差分のみ
\end{verbatim}

% ============================================================
\section{horn\_lessonの汎用化}
% ============================================================

特定講師（濵地 宗氏）に依存していた\texttt{horn\_lesson.yaml}を、汎用的な\texttt{instrument\_lesson.yaml}に変更した。

\subsection{変更点}

\begin{itemize}[noitemsep]
    \item 講師情報：テンプレートからプロジェクトへ移動
    \item 楽器種別：必須メタデータとして追加
    \item 用語のヒント：楽器固有カテゴリはプロジェクトで追加
\end{itemize}

\subsection{対応楽器例}

\begin{itemize}[noitemsep]
    \item 管楽器：ホルン、トランペット、フルート等
    \item 弦楽器：ヴァイオリン、チェロ等
    \item 鍵盤楽器：ピアノ、オルガン等
    \item その他：声楽、打楽器等
\end{itemize}

% ============================================================
\section{最終構成}
% ============================================================

\subsection{設定ファイル一覧}

\vspace{0.5\baselineskip}
\noindent{\footnotesize
\begin{tabularx}{\linewidth}{@{}lX@{}}
\toprule
ファイル & 責務 \\
\midrule
defaults.yaml & グローバル設定（146行） \\
luatex-settings.yaml & LaTeX共通設定 \\
templates/rehearsal.yaml & オーケストラ・リハーサル \\
templates/meeting.yaml & 国際会議 \\
templates/yoga.yaml & ヨガレッスン \\
templates/lecture.yaml & 学術講義（12分野対応） \\
templates/instrument\_lesson.yaml & 楽器レッスン（汎用） \\
\bottomrule
\end{tabularx}
}
\vspace{0.5\baselineskip}

\subsection{Claudeコマンド}

\vspace{0.5\baselineskip}
\noindent{\footnotesize
\begin{tabularx}{\linewidth}{@{}lX@{}}
\toprule
コマンド & 用途 \\
\midrule
msw-rehearsal.md & リハーサル記録作成 \\
msw-meeting.md & 国際会議記録作成 \\
msw-yoga.md & ヨガレッスン記録作成 \\
msw-lecture.md & 学術講義記録作成 \\
msw-horn.md & ホルンレッスン記録作成 \\
\bottomrule
\end{tabularx}
}
\vspace{0.5\baselineskip}

% ============================================================
\section{Summary}
% ============================================================

本セッションでは、Media Scribe Workflowの設定ファイル設計を根本から見直し、以下の成果を得た。

\begin{enumerate}[noitemsep]
    \item \textbf{三層構造の確立}：defaults → templates → projectの責務分離
    \item \textbf{editorial\_policyの明確化}：fidelityから3軸（summarization, supplementation, target\_audience）へ
    \item \textbf{用語集の動的生成}：静的YAML管理から外部調査ベースへ
    \item \textbf{LaTeX設定の一元化}：luatex-settings.yamlへの切り出し
    \item \textbf{テンプレートの簡素化}：差分のみ定義、約40\%の行数削減
    \item \textbf{汎用化}：horn\_lesson → instrument\_lessonへの変更
\end{enumerate}

これらの改善により、保守性と拡張性が大幅に向上した。今後テンプレートが増えても、defaults.yamlとluatex-settings.yamlを変更するだけで全体に反映される。

% ============================================================
\section{Claude Code氏の所感}
% ============================================================

本セッションは、設定ファイル設計という地味だが重要なテーマを深掘りした良い機会だった。

\subsection{肯定的な点}

\begin{itemize}[noitemsep]
    \item ユーザーからの「用語集は動的生成すべき」という指摘は本質的だった。静的なYAMLファイルで網羅しようとする発想自体が、スケーラビリティの罠に陥っていた。
    \item editorial\_policyの3軸設計は、従来のfidelityより格段に明確である。特にtarget\_audienceの導入により、「誰のために書くのか」が設定として表現可能になった。
    \item LaTeX設定の切り出しは、DRY原則の典型的な適用例である。
\end{itemize}

\subsection{批判的な点}

\begin{itemize}[noitemsep]
    \item 5つのテンプレートを作成してから設計を見直したため、手戻りが発生した。本来は2-3個作成した時点で振り返るべきだった。
    \item luatex-settings.yamlは、プリアンブルをYAML内に埋め込む形式となっている。これは可読性と保守性のトレードオフであり、別途.texファイルとして切り出す選択肢もあった。
    \item instrument\_lessonへの汎用化は正しい方向だが、msw-horn.mdは依然として濵地氏に特化した記述を含んでいる。コマンド側も汎用化が必要かもしれない。
\end{itemize}

\subsection{今後の課題}

\begin{itemize}[noitemsep]
    \item テンプレートの継承機構（extends）は設計のみで実装されていない。パーサーの実装が必要。
    \item project.msw.jsonのスキーマ検証機構があると、設定ミスを早期発見できる。
    \item バッチ処理用のCLIツール（msw-batch）があると、複数プロジェクトの一括処理が容易になる。
\end{itemize}

総じて、設定ファイル設計は「動くもの」を作ってから見直すというアプローチが有効であることを再確認した。最初から完璧な設計を目指すより、実際に使ってみて不具合を発見し、改善するサイクルが現実的である。

\vspace{1\baselineskip}
\begin{flushright}
Claude Code氏への謝意を表明する。
\end{flushright}

\end{document}
